\chapter{Fundamentação Teórica}
\label{chp:fundamentacao}

Este capítulo apresenta os conceitos que fundamentam a avaliação comparativa de
soluções de armazenamento de objetos em arquiteturas de \textit{Data Lake}
híbridas. São discutidos os conceitos de \textit{Data Lake}, armazenamento de
objetos, compatibilidade com a API S3, implantação \textit{on-premises} e em
nuvem, além dos critérios de desempenho, governança e sustentabilidade
tecnológica que orientam o estudo experimental.

\section{Arquiteturas de \textit{Data Lake}}

\textit{Data Lakes} armazenam dados em diferentes formatos e permitem que o
processamento seja realizado de forma desacoplada da camada de armazenamento.
Em arquiteturas híbridas, parte dos dados e dos serviços permanece sob controle
da organização, enquanto outra parte utiliza recursos de provedores de nuvem.

\section{Armazenamento de objetos e API S3}

O armazenamento de objetos organiza os dados como objetos associados a
metadados e identificadores, normalmente agrupados em contêineres. A API S3
tornou-se uma referência de interoperabilidade para esse modelo, permitindo que
aplicações utilizem diferentes implementações com alterações reduzidas.

\section{Critérios de avaliação}

A comparação das soluções considera desempenho, compatibilidade, operação,
licenciamento, governança e sustentabilidade. Esses critérios são detalhados e
operacionalizados a partir dos achados do mapeamento sistemático apresentado no
Capítulo~\ref{chp:mapeamento}.
